\chapter{Revisão Bibliográfica}
O estudo de comportamentos emergentes em sistemas sociais e computacionais tem suas raízes na física estatística e na teoria de grafos (Barabási, 2016). No contexto de agentes autônomos, modelos matemáticos tradicionais frequentemente assumem comportamentos estocásticos simplificados. Contudo, trabalhos recentes de Park et al. (2023) demonstraram que a integração de LLMs como motor cognitivo de agentes produz interações sociais de alta complexidade.

Diferentemente dos modelos centralizados, a inferência na borda (\textit{edge computing}) introduz restrições assíncronas e heterogeneidade cognitiva (modelos menores de 7B e 14B parâmetros). A transição de fase nestes sistemas — por exemplo, a rapidez com que a rede inteira chega à solução de um problema — requer modelagens com equações diferenciais e grafos dinâmicos para entender como o "conhecimento" se propaga e se estabiliza sem um coordenador central.
